\label{sec:results}

All numbers below come from the captured logs in \texttt{Outputs/1.txt}--\texttt{Outputs/6.txt} unless noted. Figures are generated by the experiment scripts into \texttt{graphs/} where implemented. Where we use two parallel \texttt{iperf3} flows, combined throughput is the sum of per-flow averages or per-interval sums as printed by the driver.

\subsection{Limitations}
Parallel clients bound to interfaces approximate multipath capacity; they are not always identical to a single MPTCP connection ID in the kernel. Experiment~6 studies \emph{handover-like} behavior (primary-path RTT degradation) rather than Jain-fair bandwidth sharing among independent competing flows; the latter remains a possible extension (Section~\ref{sec:future}).

\subsection{Experiment~1: TCP per-path baselines}
Table~\ref{tab:exp1} summarizes independent TCP runs on each emulated path. Throughput sits close to the configured caps (10 and 20~Mb/s), indicating sane Mininet shaping.

\begin{table}[htbp]
  \centering
  \caption{Experiment~1: baseline TCP average throughput (10~s runs).}
  \label{tab:exp1}
  \begin{tabular}{lccc}
    \toprule
    Path & Config.\ cap & Avg.\ throughput & Retransmissions \\
    \midrule
    Path~1 (Wi-Fi-like) & 10~Mb/s & 9.45~Mb/s & 0 \\
    Path~2 (LTE-like)   & 20~Mb/s & 16.95~Mb/s & 0 \\
    \bottomrule
  \end{tabular}
\end{table}

\begin{figure}[htbp]
  \centering
  \includegraphics[width=0.92\linewidth]{fig01_tcp_baseline_timeseries.png}
  \caption{Experiment~1: per-second throughput for sequential single-path TCP runs (two separate \texttt{iperf3} sessions).}
  \label{fig:exp1}
\end{figure}

\subsection{Experiment~2: TCP baselines vs.\ dual-path aggregation}
Sequential TCP baselines on the heterogeneous topology yield 9.45 and 16.99~Mb/s (combined nominal sum $26.44$~Mb/s). With two concurrent flows (one per interface), we observe per-path averages of 9.49 and 17.67~Mb/s and a reported combined average of 27.16~Mb/s, i.e.\ $\approx$102.7\% ``aggregation efficiency'' relative to the sum of independent TCP measurements---consistent with statistical variation and Mininet timing rather than a protocol violation. Figure~\ref{fig:exp2} contrasts per-path TCP time series with the parallel-stream run.

\begin{figure}[htbp]
  \centering
  \begin{subfigure}[b]{0.92\linewidth}
    \centering
    \includegraphics[width=\linewidth]{fig02_tcp_baseline_timeseries.png}
    \caption{Sequential TCP baselines (one path at a time).}
  \end{subfigure}

  \vspace{0.6em}
  \begin{subfigure}[b]{0.92\linewidth}
    \centering
    \includegraphics[width=\linewidth]{fig02_mptcp_parallel_streams_timeseries.png}
    \caption{Parallel \texttt{iperf3} streams, one bound to each interface (MPTCP endpoints enabled).}
  \end{subfigure}
  \caption{Experiment~2: throughput over time for baselines vs.\ dual-path aggregation.}
  \label{fig:exp2}
\end{figure}

\subsection{Experiment~3: Sudden path failure}
At $t=10$~s the receiver's first-path interface is taken down and the corresponding client is stopped. The LTE-side stream continues for the full 30~s window with average throughput 18.37~Mb/s and zero retransmissions on that stream in our log, while the Wi-Fi stream shows throughput dropping to zero after the failure. This supports \textbf{session continuity} in the sense that the application-level transfer on the surviving path need not be restarted; aggregate goodput, however, falls to single-path levels after the event. Figure~\ref{fig:exp3} marks the failure instant on the combined throughput trace.

\begin{figure}[htbp]
  \centering
  \includegraphics[width=0.92\linewidth]{fig03_path_failure_throughput.png}
  \caption{Experiment~3: combined throughput with vertical line at path failure ($t=10$~s).}
  \label{fig:exp3}
\end{figure}

\subsection{Experiment~4: Bandwidth collapse on one path}
Path~1 and Path~2 each start at 10~Mb/s and $\approx$40~ms RTT. Dual-stream baselines achieve combined $\approx$18.76~Mb/s, slightly above the sum of independent TCP baselines (18.56~Mb/s), matching the printed 101.1\% aggregation efficiency. After Path~2 is throttled to 2~Mb/s at $t=10$~s, the driver's windowed averages report combined throughput falling from $\approx$24.51~Mb/s before collapse to $\approx$15.33~Mb/s after, i.e.\ \textbf{62.5\% retained}. Path~1's average rises from $\approx$12.25 to $\approx$13.71~Mb/s in the post-collapse window, indicating load shifting toward the stable path; Path~2 shows many retransmissions (1863 over 25~s in our log) as congestion control reacts to the bottleneck. Figure~\ref{fig:exp4} contrasts the pre-collapse aggregation shape with the collapse scenario.

\begin{figure}[htbp]
  \centering
  \begin{subfigure}[b]{0.92\linewidth}
    \centering
    \includegraphics[width=\linewidth]{fig04_mptcp_baseline_timeseries.png}
    \caption{Dual-path baseline before any impairment.}
  \end{subfigure}

  \vspace{0.6em}
  \begin{subfigure}[b]{0.92\linewidth}
    \centering
    \includegraphics[width=\linewidth]{fig04_bandwidth_collapse_timeseries.png}
    \caption{Dual-path run with Path~2 throttled to 2~Mb/s at $t=10$~s (marker in plot).}
  \end{subfigure}
  \caption{Experiment~4: symmetric topology---baseline vs.\ bandwidth collapse on Path~2.}
  \label{fig:exp4}
\end{figure}

\subsection{Experiment~5: RTT spike on Path~2}
Injecting $\approx$800~ms RTT on Path~2 at $t=10$~s while using TCP \emph{only} on Path~2 reduces average throughput for that run to 4.22~Mb/s, with before/after spike window averages of 11.93 and 1.61~Mb/s (13.5\% throughput retained on that single path). Under dual-path transfer, combined throughput retained across the same event is reported as \textbf{71.0\%} (23.85~Mb/s before vs.\ 16.94~Mb/s after), with Path~2's post-spike average dropping to 4.03~Mb/s while Path~1 increases modestly (11.92 to 12.91~Mb/s). Figure~\ref{fig:exp5a} shows the multipath time series; Figure~\ref{fig:exp5b} summarizes the TCP vs.\ MPTCP comparison from the driver's final table.

\begin{figure}[htbp]
  \centering
  \includegraphics[width=0.92\linewidth]{fig05_mptcp_rtt_spike_timeseries.png}
  \caption{Experiment~5: per-path and combined throughput during Path~2 RTT spike (dual streams).}
  \label{fig:exp5a}
\end{figure}

\begin{figure}[htbp]
  \centering
  \includegraphics[width=0.75\linewidth]{fig05_summary_post_spike_comparison.png}
  \caption{Experiment~5: summary comparison (throughput retained under spike; TCP single-path vs.\ multipath-style aggregate).}
  \label{fig:exp5b}
\end{figure}

\begin{table}[htbp]
  \centering
  \caption{Experiment~5: selected metrics from the automated summary table (post-spike windows).}
  \label{tab:exp5}
  \begin{tabular}{lc}
    \toprule
    Metric & Value \\
    \midrule
    TCP Path~2 only --- avg.\ during full 25~s run & 4.22~Mb/s \\
    TCP Path~2 only --- throughput retained (post-spike window) & 13.5\% \\
    MPTCP-style dual stream --- combined throughput retained & 71.0\% \\
    Ratio: MPTCP combined avg.\ vs.\ TCP Path~2 only (driver table) & $\approx 10.5{\times}$ \\
    \bottomrule
  \end{tabular}
\end{table}

\subsection{Experiment~6: Handover stress (primary-path RTT ramp)}
\label{subsec:exp6}

This run (\texttt{6.py}, log \texttt{Outputs/6.txt}) uses two $10$~Mb/s paths with Path~1 initially lower RTT ($20{+}20$~ms per link on \texttt{TCLink}) and Path~2 higher RTT ($50{+}50$~ms), analogous to a ``Wi-Fi'' vs.\ ``LTE'' preference. Kernel MPTCP uses the \texttt{blest} packet scheduler. Part~1--2 repeat TCP per-path baselines and a $15$~s dual-stream aggregation baseline. Part~3 runs parallel streams for $25$~s; after $t=10$~s, Path~1 one-way delay is stepped upward via \texttt{tc netem} on host interfaces and \texttt{s1} ports (target RTT milestones $100$--$250$~ms) so the stack should shift traffic toward Path~2. \textbf{No matplotlib figure} was produced for this script; we rely on the driver's summary and the per-second table below.

Table~\ref{tab:exp6baseline} gives baselines; Table~\ref{tab:exp6handover} summarizes the handover phase. Combined average throughput changes only slightly across the event boundary ($24.89$~Mb/s before vs.\ $23.39$~Mb/s after), while Path~1's average in those windows falls from $11.94$ to $10.49$~Mb/s and Path~2's post-handover average is $12.91$~Mb/s---consistent with modest load migration. Path~1 accumulated $493$ retransmissions over the $25$~s handover run (Path~2: $0$ in our log), reflecting congestion dynamics on the degraded path. The per-second trace shows a short dip at $t{=}12$~s (Path~1 reported $0$~Mb/s, total $10.49$~Mb/s) before recovery, illustrating transient rebalancing rather than a hard failure.

\begin{table}[htbp]
  \centering
  \caption{Experiment~6: TCP baselines and dual-stream aggregation (from \texttt{Outputs/6.txt}).}
  \label{tab:exp6baseline}
  \begin{tabular}{lcc}
    \toprule
    Phase & Quantity & Value \\
    \midrule
    \multirow{3}{*}{TCP (10~s each)} & Path~1 avg & 9.28~Mb/s \\
     & Path~2 avg & 8.68~Mb/s \\
     & Sum (theoretical cap $\sim$20~Mb/s) & 17.97~Mb/s \\
    \midrule
    \multirow{4}{*}{Dual-stream baseline (15~s)} & Path~1 avg & 9.38~Mb/s \\
     & Path~2 avg & 8.98~Mb/s \\
     & Combined avg & 18.36~Mb/s \\
     & Aggregation efficiency vs.\ TCP sum & 102.2\% \\
    \bottomrule
  \end{tabular}
\end{table}

\begin{table}[htbp]
  \centering
  \caption{Experiment~6: handover run ($25$~s; Path~1 RTT stepped after $t{=}10$~s). Combined averages are windowed by the driver.}
  \label{tab:exp6handover}
  \begin{tabular}{lc}
    \toprule
    Metric & Value \\
    \midrule
    Combined throughput avg.\ \emph{before} handover window & 24.89~Mb/s \\
    Combined throughput avg.\ \emph{after} handover window & 23.39~Mb/s \\
    Path~1 avg before / after (driver windows) & 11.94 / 10.49~Mb/s \\
    Path~2 avg after handover (driver) & 12.91~Mb/s \\
    Path~1 / Path~2 avg over full $25$~s run & 9.27 / 9.20~Mb/s \\
    Retransmissions (Path~1 / Path~2) & 493 / 0 \\
    \bottomrule
  \end{tabular}
\end{table}

\begin{table}[htbp]
  \centering
  \caption{Experiment~6: excerpt of per-second combined throughput (Mb/s) around the handover (from log).}
  \label{tab:exp6timeseries}
  \begin{tabular}{crrr}
    \toprule
    $t$ (s) & Path~1 & Path~2 & Total \\
    \midrule
    8 & 10.48 & 10.49 & 20.97 \\
    9 & 10.49 & 10.49 & 20.97 \\
    10 & 10.48 & 10.48 & 20.97 \\
    11 & 10.49 & 10.49 & 20.98 \\
    12 & 0.00 & 10.49 & 10.49 \\
    13 & 10.49 & 20.99 & 31.48 \\
    14 & 10.49 & 10.48 & 20.97 \\
    15 & 10.49 & 10.48 & 20.97 \\
    \bottomrule
  \end{tabular}
\end{table}

Overall, path failure and RTT-spike experiments show the qualitative benefit of maintaining a second usable path: aggregate throughput degrades but the transfer is not confined to the impaired link alone. Experiment~6 complements this by showing that gradual primary-path degradation can be absorbed with limited aggregate loss, at the cost of retransmissions on the degrading path.
